\section{Introduction}\label{sec:intro}

% Register (PI, revision of draft 91): the point and its evidence in one movement, definitions in
% appositions, no announcing sentences, no "asks whether", one figure reference (the design), no
% findings by figure or table number. P1 is the PI's model paragraph with the facts checked.
Foundation models~\cite{bommasani}, networks pretrained once on a large dataset and then adapted to
many tasks, have become standard in language~\cite{bert,gptthree} and vision~\cite{mahajan,mae}, where
they adapt to new tasks from few labeled examples, and the same approach now predicts protein
structure from sequence~\cite{esmfold} and analyzes whole collision events~\cite{evenet}. What such a
model carries to a new task is fixed during pretraining, by the data it sees and by what it is trained
to predict, and jet physics is an unusually clean place to study the second of these. The experiments
at the Large Hadron Collider (LHC)~\cite{lhc} reconstruct every proton--proton collision as a set of
particles; simulation supplies labeled collisions in essentially unlimited numbers; and each simulated
jet carries the exact decay that produced it, so the same jets can be labeled as finely or as coarsely
as one chooses.

% P2. Jets, tagging, jet foundation models; supervised models serve many analyses through summed
% classes, and their label sets grow.
A top quark, a $W$, $Z$ or Higgs boson, or a new resonance produced with momentum well above its mass
decays into quarks or leptons whose showers merge into a single jet, a collimated spray of particles,
and analyses identify the particle from the structure of that jet~\cite{larkoski,kogler,cmshbb,cmszprime}:
its number of prongs, the separate energy clusters left by the decay products; the flavor of its quarks,
which a $b$ or $c$ hadron reveals through a displaced vertex a few millimeters from the collision
point~\cite[Sec.~4.1]{cmsbtag}; and its mass. The taggers that do this are deep networks that read
every particle in the jet, trained on simulated jets to separate the signal decays from jets of light
quarks and gluons, the quantum chromodynamics (QCD) background~\cite{particlenet,part,gntwo,cmspart}.
Jet foundation models are pretrained either without labels, by reconstructing masked particles or
generating jets~\cite{mpm,omnijet,hepjepa}, or as classifiers of simulated
jets~\cite{part,omnilearn,sophon}, the reference objective against which the self-supervised ones are
measured~\cite{objectives}. A classifier of this kind separates many fine classes and serves many
analyses, each of which sums the fine classes into the signal and background it
needs~\cite{sophon,glopartdp}, and its label set, the classes it is trained to separate, keeps growing:
the Sophon model is pretrained on the \Kfull{} classes of the JetClass-II
dataset~\cite{sophon,jetclasstwo}, and the CMS tagger GloParT went from \LitGlopartVoneNClasses{}
classes in its first version to \LitGlopartVthreeNClasses{} in its third~\cite{glopartdp}.

% P3. The gap: Sophon's properties hold at the output of an ideal classifier; the backbone is
% untested because every published comparison changed the jets with the labels.
At the output of an ideal classifier finer labels cost nothing, since the score of a class split into
exclusive subclasses equals the sum of the subclass scores and adding classes leaves the ratios between
the scores of the original classes unchanged~\cite[Sec.~III, Properties~1 and~2]{sophon}; a model
trained on fine classes therefore serves every coarser definition of the signal. These properties are
silent about the backbone, the network below the classification head, which is the part reused when
the model is adapted to a new task, and whether the label set shapes it has not been tested, because
every published comparison changed the training jets together with the labels. Sophon's
\LitSophonAblationNClasses-class model, which identifies $X\to b\bar b$ decays less well than the
\Kfull-class model, was trained only on two-prong decays and QCD jets and on fewer jets per
epoch~\cite[Sec.~III, App.~B.2]{sophon}. In the thesis that introduced GloParT, two-prong and
top-quark classes added to a classifier of $H\to WW^*$ decays came with their own training jets and
raised the area under the receiver operating characteristic curve (AUC) against QCD jets from
\LitGlopartAddedAucBase{} to \LitGlopartAddedAucTop, a change the author reads as
minor~\cite[Sec.~11.4.1]{glopart}.\footnote{Studies that vary the composition of the pretraining data
also change the labels together with the jets: Uslu et~al.~\cite{datacomp} change which JetClass-II
classes enter pretraining at a fixed total number of jets, and SuperJEPA~\cite{superjepa} pretrains on
a subset of the JetClass classes and on all of them with different numbers of jets.}

% P4. Vision: the answer depends on the grouping and on how the features are used; Mahajan's
% probe-versus-fine-tuning result is the one that motivates our design.
Image classification has studied the same choice, and its answer depends on how the classes are
grouped and how the pretrained network is used. Coarse groups of the ImageNet classes cost little in
detection after fine-tuning~\cite{huh}, and finer labels that follow a taxonomy of species help a
fine-tuned network up to a point, whereas class identities assigned at random to individual images do
worse than coarse labels at every granularity~\cite{hong}. A larger label set of web hashtags gives a
clear lead to a linear probe, a logistic regression trained on the frozen features (the backbone's
output with its weights held fixed), but the lead nearly disappears under fine-tuning, which continues
training all weights on the new task~\cite{mahajan}; jets have not been studied this way.

% P5. Design, with the one figure reference; the hierarchy in one sentence (Methods has the rest).
We pretrain Sophon's network on the same JetClass-II jets, in the same order and from the same initial
weights, and change only the label set and with it the width of the output layer, as
figure~\ref{fig:design} shows. The label sets are four nested levels of a class hierarchy, each
pretrained \MlstNRunsWord{} times, in which every merge removes a named set of distinctions, from the
flavor content of QCD jets at the first step to quark flavor itself at the last, leaving \Kcoarse{}
classes that keep only the number and type of the prongs. Two kinds of control hold the class count at
\Kcoarse: random partitions, which assign the resonant classes to groups at random, and partitions that
differ only in a flavor boundary, a class boundary between decays that differ only in quark flavor. We
read the frozen features with linear probes and with MLP probes, small multilayer perceptrons, and
fine-tune the pretrained models on JetClass-II, on the JetClass dataset~\cite{part}, on top
tagging~\cite{toptag} and on quark/gluon discrimination with Pythia and Herwig
samples~\cite{efn,qgpythia,qgherwig}, with Sophon's $X\to b\bar b$ and $X\to bs$ benchmark as the % CHECK: pending Xbx* (drop this clause with section 4.7 if the samples are not staged by 12 Nov)
last~\cite{sophon}.

\begin{figure}[tbp]
\centering
\paperfig{\linewidth}{F1}
\caption{Design of the study. (a)~One network is pretrained on the same jets four times, once for each of
four nested label sets with \Kfull, \Kres, \Kmid{} and \Kcoarse{} classes, through a hidden layer of
\DesignHeadHidden{} units; each merged class spans the rows it merges, rows are schematic, and below each
merge is the distinction it removes. (b)~The pretrained network is then used in two ways: its frozen
features, the output of the backbone with its weights held fixed, feed a linear probe, or a new head with
its own hidden layer of \DesignHeadHidden{} units is added and all weights are fine-tuned. The colors of
the label sets are those of every later figure.}
\label{fig:design}
\end{figure}

% P6. Findings in one paragraph, no figure or table numbers, no numbers.
% CHECK: MlstCmpJcTwoResResOverFullPr reads "equal within" and MlstCmpJcTwoQcdResOverFull > 1 by more than its SD
% CHECK: MlstCmpJcTwoCoarseOverMid > MlstCmpJcTwoMidOverRes > MlstCmpJcTwoResOverFull; MlstPerClassCoarseTopClass is an e/mu or tau_e/tau_mu class; MlstPerClassCoarseQuarkShare > MlstPerClassCoarseLeptonShare
% CHECK: MlstCmpJcTwoRandOverCoarse < 1; MlstCmpJcTwoF{one,oner}OverFzero < 1 (both runs)
% CHECK: MlstFtJcOneCoarseNRejDrop >= two and MlstFtJcOne{Res,Mid}NRejDrop = none, on the long-budget cells once all 30 are in
% CHECK: MlstFtTopRthirty{Res,Mid,Coarse}OverFull and MlstFtQgRthirty{Res,Mid,Coarse}OverFull within a few per cent of one
% CHECK: MlstLpTopRthirtyFullOverInitBnZeroPr reads "equal within" and MlstLpTopRthirtyDrawOverZeroFull well above one
% CHECK: MlstCmpMlpResMassOverRes < 1 and MlstSculptCoarseMassOverCoarseOne < 1; MlstCmpJcTwoResMassOverRes within a few per cent of one
The frozen features keep what the labels ask the network to separate: merging the QCD classes costs only
on those classes; summed over classes, the merge of quark flavors is the costliest step of the
hierarchy, while class by class the largest losses fall on decays that differ only in lepton flavor; and
at a fixed number of classes the random partitions lose less than the hierarchy, so the cost depends on
which classes are merged and not only on how many remain, with any flavor boundary keeping flavor
readable. Fine-tuning narrows these differences without removing them: on JetClass, a separate
simulation with every input present, the \Kcoarse-class labels still lower the background rejection for
several signal classes. The public top-tagging and quark/gluon benchmarks lack some of the network's
inputs, and on them the label sets differ far less; with the missing inputs set to zero, a linear probe
on the pretrained features does no better than one on a randomly initialized network, and on top
tagging replacing the zeros with values a real particle takes restores much of the difference. An auxiliary regression of the
jet mass, trained alongside the classification, makes the frozen features carry the mass and reduces
the sculpting of the background mass distribution, at a small cost to the classification.

The rest of this paper is organized as follows. Section~\ref{sec:related} places the study among work
on pretraining objectives, label granularity, auxiliary targets and transfer; section~\ref{sec:methods}
describes the data, the label sets and the evaluation; section~\ref{sec:results} gives the results task
by task; and section~\ref{sec:discussion} discusses them and states the limitations.
